\documentclass{article}

\usepackage{ltablex}
\usepackage{booktabs}
\usepackage{pdflscape}

\keepXColumns

\newcolumntype{Y}{>{\raggedright\arraybackslash}X}
\newcolumntype{P}[1]{>{\raggedright\arraybackslash}p{#1}}

\title{Reflection and Traceability Report on \progname}

\author{\authname}

\date{}

\input{../Comments}
%% Common Parts

\newcommand{\progname}{Software Engineering} % PUT YOUR PROGRAM NAME HERE
\newcommand{\authname}{Team 8, AlgoCatan
\\ Jake Read
\\ Rebecca Di Filippo
\\ Matthew Cheung
\\ Sunny Yao} % AUTHOR NAMES

\usepackage{hyperref}
    \hypersetup{colorlinks=true, linkcolor=blue, citecolor=blue, filecolor=blue,
                urlcolor=blue, unicode=false}
    \urlstyle{same}
                                


\begin{document}

\maketitle

\plt{Reflection is an important component of getting the full benefits from a
learning experience.  Besides the intrinsic benefits of reflection, this
document will be used to help the TAs grade how well your team responded to
feedback.  Therefore, traceability between Revision 0 and Revision 1 is and
important part of the reflection exercise.  In addition, several CEAB (Canadian
Engineering Accreditation Board) Learning Outcomes (LOs) will be assessed based
on your reflections.}

\section{Changes in Response to Feedback}

\plt{Summarize the changes made over the course of the project in response to
feedback from TAs, the instructor, teammates, other teams, the project
supervisor (if present), and from user testers.}

\plt{For those teams with an external supervisor, please highlight how the feedback 
from the supervisor shaped your project.  In particular, you should highlight the 
supervisor's response to your Rev 0 demonstration to them.}

\plt{Version control can make the summary relatively easy, if you used issues
and meaningful commits.  If you feedback is in an issue, and you responded in
the issue tracker, you can point to the issue as part of explaining your
changes.  If addressing the issue required changes to code or documentation, you
can point to the specific commit that made the changes.  Although the links are
helpful for the details, you should include a label for each item of feedback so
that the reader has an idea of what each item is about without the need to click
on everything to find out.}

\plt{If you were not organized with your commits, traceability between feedback
and commits will not be feasible to capture after the fact.  You will instead
need to spend time writing down a summary of the changes made in response to
each item of feedback.}

\plt{You should address EVERY item of feedback.  A table or itemized list is
recommended.  You should record every item of feedback, along with the source of
that feedback and the change you made in response to that feedback.  The
response can be a change to your documentation, code, or development process.
The response can also be the reason why no changes were made in response to the
feedback.  To make this information manageable, you will record the feedback and
response separately for each deliverable in the sections that follow.}

\plt{If the feedback is general or incomplete, the TA (or instructor) will not
be able to grade your response to feedback.  In that case your grade on this
document, and likely the Revision 1 versions of the other documents will be
low.} 

\newpage
\begin{landscape}
\subsection{SRS and Hazard Analysis}

\begin{tabularx}{\linewidth}{|P{2.3cm}|Y|P{2.5cm}|P{3.5cm}|}

        % Header
        \caption{SRS and HA Feedback and Resultant Changes} \label{TblSRSHA} \\
        \hline
        \textbf{Source} &
        \textbf{Feedback Summary} &
        \textbf{Link to Issue} &
        \textbf{Link to Commit} \\
        \hline
        \endfirsthead

        % Subsequent pages header
        \hline
        \textbf{Source} &
        \textbf{Feedback Summary} &
        \textbf{Link to Issue} &
        \textbf{Link to Commit} \\
        \hline
        \endhead

        % Prompt that the table continues
        \hline
        \multicolumn{4}{r}{\textit{Continued on next page}} \\
        \endfoot

        % No footer on last page
        \hline
        \endlastfoot

        % Table Content
        Peer Review &
        NFR-E.10 and 11 missing values &
        \href{https://github.com/AlgoCatan/RLCatan/issues/16}{Issue \#16} \newline \href{https://github.com/AlgoCatan/RLCatan/issues/17}{Issue \#17} &
        \href{https://github.com/AlgoCatan/RLCatan/commit/983904c2f1b365fc0b136a427e5a16b1c1ac575e}{Commit 983904c} \\

        \hline
        Peer Review &
        Some stakeholders not defined &
        \href{https://github.com/AlgoCatan/RLCatan/issues/13}{Issue \#13} &
        \href{https://github.com/AlgoCatan/RLCatan/commit/d21fdf5aec305c48e0995328094454a36437e8d9}{Commit d21fdf5} \\

        \hline
        Peer Review &
        Section S.3 missing important APIs/interfaces &
        \href{https://github.com/AlgoCatan/RLCatan/issues/19}{Issue \#19} &
        \href{https://github.com/AlgoCatan/RLCatan/commit/bd9db5308f045a4b8a9615be7ca28d0ee1123aa3}{Commit bd9db53} \\

        \hline
        Peer Review &
        Invariants appear to be duplicated &
        \href{https://github.com/AlgoCatan/RLCatan/issues/18}{Issue \#18} &
        \href{https://github.com/AlgoCatan/RLCatan/commit/4574535067f674fb58bea087aebc4723ab7042a7}{Commit 4574535} \\

        \hline
        Peer Review &
        Not all high-level usage scenarios given later details &
        \href{https://github.com/AlgoCatan/RLCatan/issues/14}{Issue \#14} &
        \href{https://github.com/AlgoCatan/RLCatan/commit/849db38b2cf12eb7b75cb533febe9b18bc03d4c8}{Commit 849db38} \\

        \hline
        Peer Review &
        Detailed Usage Scenario poorly detailed and justification from dev rather than user perspective &
        \href{https://github.com/AlgoCatan/RLCatan/issues/15}{Issue \#15} &
        \href{https://github.com/AlgoCatan/RLCatan/commit/849db38b2cf12eb7b75cb533febe9b18bc03d4c8}{Commit 849db38} \\

        \hline
        Peer Review &
        Too few details in high-level usage scenarios &
        \href{https://github.com/AlgoCatan/RLCatan/issues/20}{Issue \#20} &
        N/A (See issue) \\

        \hline
        Peer Review &
        Unclear invariant source &
        \href{https://github.com/AlgoCatan/RLCatan/issues/21}{Issue \#21} &
        \href{https://github.com/AlgoCatan/RLCatan/commit/77bdeaca970d2643ddce250b1b030626451bfc37}{Commit 77bdeac} \\

        \hline
        Peer Review &
        Add NFRs specific to each component &
        \href{https://github.com/AlgoCatan/RLCatan/issues/23}{Issue \#23} &
        N/A (See issue) \\

        \hline
        Peer Review &
        HA missing legal move assumption &
        \href{https://github.com/AlgoCatan/RLCatan/issues/24}{Issue \#24} &
        N/A (See issue) \\

        \hline
        Peer Review &
        HA camera safety req has too few implementation details &
        \href{https://github.com/AlgoCatan/RLCatan/issues/26}{Issue \#26} &
        N/A (See issue) \\

        \hline
        Peer Review &
        Missing potential loss incurred by system failure &
        \href{https://github.com/AlgoCatan/RLCatan/issues/27}{Issue \#27} &
        \href{https://github.com/AlgoCatan/RLCatan/commit/c1f5ec82a0a6f454403809555ce19ad0a5551874}{Commit c1f5ec8} \\

        \hline
        Peer Review &
        FR.Sa.11 not verifiable &
        \href{https://github.com/AlgoCatan/RLCatan/issues/28}{Issue \#28} &
        \href{https://github.com/AlgoCatan/RLCatan/commit/5258f02568fac29461190a2092b85c8cf5e54813}{Commit 5258f02} \\

        \hline
        Peer Review &
        FMEA table missing fields &
        \href{https://github.com/AlgoCatan/RLCatan/issues/29}{Issue \#29} &
        N/A (See issue) \\

        \hline
        TA Feedback &
        The wording of the requirements was very poor, and did not follow the typical pattern of "The system shall...". (This was one of the major areas we lost grades on the SRS doc) &
        \href{https://github.com/AlgoCatan/RLCatan/issues/109}{Issue \#109} &
        \href{https://github.com/AlgoCatan/RLCatan/commit/c3c3f023ea6e9fe4c5fcb4556fc0c172765f1e45}{Commit c3c3f02} \\

        \hline
        Peer Review &
        Section S.1 needed clearer component interactions and links between components &
        \href{https://github.com/AlgoCatan/RLCatan/issues/22}{Issue \#22} &
        \href{https://github.com/AlgoCatan/RLCatan/commit/57547fead57e4b5eda825a54963c6f3978c4d145}{Commit 57547fe} \\

        \hline
        Peer Review &
        SRS document still has wording issues &
        \href{https://github.com/AlgoCatan/RLCatan/issues/22}{Issue \#22} &
        \href{https://github.com/AlgoCatan/RLCatan/commit/57547fead57e4b5eda825a54963c6f3978c4d145}{Commit 57547fe} \\

        \hline
        Peer Review &
        MIS M8 resetQueue() should not be local-only; promoted to exported access routine &
        \href{https://github.com/AlgoCatan/RLCatan/issues/49}{Issue \#49} &
        \href{https://github.com/AlgoCatan/RLCatan/commit/769daf0a1f2ce1a065e57fa003b05a77939abfba}{Commit 769daf0} \\

        \hline
        Peer Review &
        MIS integer constants/state variables (e.g., max turns, score) needed tighter numeric domains &
        \href{https://github.com/AlgoCatan/RLCatan/issues/51}{Issue \#51} &
        \href{https://github.com/AlgoCatan/RLCatan/commit/64a9e404d4ee43724bc3a08603f72db58f78df77}{Commit 64a9e40} \\

\end{tabularx}

\newpage
\subsection{Design and Design Documentation}
We did not receive any peer reviews for the MG or MIS documents, hence the shorter feedback table for this section.

\begin{tabularx}{\linewidth}{|P{2.3cm}|Y|P{2.5cm}|P{3.5cm}|}

        % Header
        \caption{Design and Design Documentation Feedback and Resultant Changes} \label{TblDesign} \\
        \hline
        \textbf{Source} &
        \textbf{Feedback Summary} &
        \textbf{Link to Issue} &
        \textbf{Link to Commit} \\
        \hline
        \endfirsthead

        % Subsequent pages header
        \hline
        \textbf{Source} &
        \textbf{Feedback Summary} &
        \textbf{Link to Issue} &
        \textbf{Link to Commit} \\
        \hline
        \endhead

        % Prompt that the table continues
        \hline
        \multicolumn{4}{r}{\textit{Continued on next page}} \\
        \endfoot

        % No footer on last page
        \hline
        \endlastfoot

\end{tabularx}

\newpage
\subsection{VnV Plan and Report}

\begin{tabularx}{\linewidth}{|P{2.3cm}|Y|P{2.5cm}|P{3.5cm}|}

        % Header
        \caption{VnV Plan and Report Feedback and Resultant Changes} \label{TblVnV} \\
        \hline
        \textbf{Source} &
        \textbf{Feedback Summary} &
        \textbf{Link to Issue} &
        \textbf{Link to Commit} \\
        \hline
        \endfirsthead

        % Subsequent pages header
        \hline
        \textbf{Source} &
        \textbf{Feedback Summary} &
        \textbf{Link to Issue} &
        \textbf{Link to Commit} \\
        \hline
        \endhead

        % Prompt that the table continues
        \hline
        \multicolumn{4}{r}{\textit{Continued on next page}} \\
        \endfoot

        % No footer on last page
        \hline
        \endlastfoot

        % Table Content
        Peer Review &
        Potential use of Postman for API testing &
        \href{https://github.com/AlgoCatan/RLCatan/issues/30}{Issue \#30} &
        N/A (See issue) \\

        \hline
        Peer Review &
        SRS peer review plan needs more details &
        \href{https://github.com/AlgoCatan/RLCatan/issues/31}{Issue \#31} &
        \href{https://github.com/AlgoCatan/RLCatan/commit/05dd5ddb17dbe99b8b6f099b12636d3b7e3cb891}{Commit 05dd5dd} \\

        \hline
        Peer Review &
        Duplicate heading in VnV plan &
        \href{https://github.com/AlgoCatan/RLCatan/issues/32}{Issue \#32} &
        \href{https://github.com/AlgoCatan/RLCatan/commit/2a4d0c932bf8fddf9a932078f75677b8e6160d89}{Commit 2a4d0c} \\

        \hline
        Peer Review &
        Section 3.7 needed more explicit validation metrics for supervisor demo and usability survey &
        \href{https://github.com/AlgoCatan/RLCatan/issues/34}{Issue \#34} &
        \href{https://github.com/AlgoCatan/RLCatan/commit/d928005ca9b6ee51cdad6ded3c3f954f5fed8d84}{Commit d928005} \\

        \hline
        Peer Review &
        Functional requirement traceability table was unclear relative to the non-functional requirements section &
        \href{https://github.com/AlgoCatan/RLCatan/issues/35}{Issue \#35} &
        \href{https://github.com/AlgoCatan/RLCatan/commit/069c83ef1850ceb88d07440ca65f05e14de7216b}{Commit 069c83e} \\

        \hline
        Peer Review &
        Module Guide section 7 needed clearer "Implemented By" technology details &
        \href{https://github.com/AlgoCatan/RLCatan/issues/48}{Issue \#48} &
        \href{https://github.com/AlgoCatan/RLCatan/commit/42ad9fb5b28da8d6b75f4038f4a749cb29cb86bf}{Commit 42ad9fb} \\

        \hline
        Peer Review &
        Module Guide UC.1 needed clearer mapping to stable software decision/rules dependency &
        \href{https://github.com/AlgoCatan/RLCatan/issues/52}{Issue \#52} &
        \href{https://github.com/AlgoCatan/RLCatan/commit/0b0977e662e96894e9aef643d61510918c1789a9}{Commit 0b0977e} \\

        \hline
        Peer Review &
        Module Guide section 9 use hierarchy diagram should clarify module dependencies and acknowledge cycles from decoupled architecture &
        \href{https://github.com/AlgoCatan/RLCatan/issues/53}{Issue \#53} &
        \href{https://github.com/AlgoCatan/RLCatan/commit/4a4f580d8d8217572d8f534f5bf5d651b84a5b51}{Commit 4a4f580} \\

        \hline
        Peer Review &
        Module Guide section 6 Clearer Connection between Requirements and Design &
        \href{https://github.com/AlgoCatan/RLCatan/issues/54}{Issue \#54} &
        \href{https://github.com/AlgoCatan/RLCatan/commit/4aa70a05f56cabdd9d2b5ad0f77149d595291913}{Commit 4aa70a0} \\

        \hline
        Peer Review &
        VnV Report section 7 needed clearer priority/scope categorization for testing-driven changes &
        \href{https://github.com/AlgoCatan/RLCatan/issues/81}{Issue \#81} &
        \href{https://github.com/AlgoCatan/RLCatan/commit/20306b9df088b1c48930f2e38f7cfb656a0afc7b}{Commit 20306b9} \\

        \hline
        Peer Review &
        VnV Report and Plan needed more specificity on bash scripts used and rationale for unit testing framework choices &
        \href{https://github.com/AlgoCatan/RLCatan/issues/82}{Issue \#82} &
        \href{https://github.com/AlgoCatan/RLCatan/commit/3fe0c62a9cc5c39628fe42c1999d10716904911b}{Commit 3fe0c62} \\

        \hline
        Peer Review &
        VnV Report usability section needed clearer explanation of test categories and expansion on Catan-to-Bot intuition &
        \href{https://github.com/AlgoCatan/RLCatan/issues/83}{Issue \#83} &
        \href{https://github.com/AlgoCatan/RLCatan/commit/9d6e0be969beadc8aed7f09c2e0f439f7a3b094d}{Commit 9d6e0be} \\

        \hline
        Peer Review &
        VnV Report Performance test needed specific timing metrics to validate performance requirements &
        \href{https://github.com/AlgoCatan/RLCatan/issues/84}{Issue \#84} &
        \href{https://github.com/AlgoCatan/RLCatan/commit/3b447b1599b62d4f24cad65d1ad63126a9243473}{Commit 3b447b1} \\

        \hline
        Peer Review &
        VnV Report NFR evaluation sections needed more detailed testing methodology and process descriptions &
        \href{https://github.com/AlgoCatan/RLCatan/issues/85}{Issue \#85} &
        \href{https://github.com/AlgoCatan/RLCatan/commit/e7d079bb45b87278a19a4e9f3b7cd4f4b22afd91}{Commit e7d079b} \\

        \hline
        Peer Review &
        VnV Report Data Integrity tests needed additional unit/component tests for invalid state and action validation &
        \href{https://github.com/AlgoCatan/RLCatan/issues/86}{Issue \#86} &
        \href{https://github.com/AlgoCatan/RLCatan/commit/2ac373a4f7b5ae4aef476faa8db9c3906013a38c}{Commit 2ac373a} \\

        \hline
        Peer Review &
        VnV Plan Performance and Scalability section needed specific performance thresholds and separate test cases &
        \href{https://github.com/AlgoCatan/RLCatan/issues/33}{Issue \#33} &
        \href{https://github.com/AlgoCatan/RLCatan/commit/47af5ac567f323701dddd3de27dd0fc482817b64}{Commit 47af5ac} \\



\end{tabularx}

\end{landscape}

\section{Extras}

\plt{Summarize the extras that were tackled by this project.  Extras
can include usability testing, code walkthroughs, user documentation, formal
proof, GenderMag personas, Design Thinking, etc.  Extras should have already
been approved by the course instructor as included in your problem statement.}

For our project, we chose to conduct usability testing, and create a user instructional video for our extras.

The usability testing was conducted in multiple stages throughout the process.
We ran the first interviews and subsequent survey following our rev0 demonstration, and the feedback we
received heavily guided our design decisions between rev0 and rev1.
As such, we will expand further on the details of this feedback in the following design iteration section.
Following the incorporation of feedback from this first round of usability testing, we ran a second round
of interviews and surveys to evaluate the changes we made in response to the first round of feedback.
We saw a noticeable increase in the UX metrics we were tracking, indicating our changes were valuable.
For more information on the specifics of this process, see the usability testing document in the extras folder.

The user instructional video was created following the completion of our rev1 documentation
and is intended to be a resource for users to learn how to use our product.
It's around 3 minutes long, and gives explanations of the various settings and features present in our
front-end interface, as well as a demonstration of how to use the product to play a game of Catan against the bot.
The video is available on our YouTube channel, and the link to the video is included in the extras folder.
The responses from users we showed it to are also listed in a doc in the same folder.

\section{Design Iteration (LO11 (PrototypeIterate))}

\plt{Explain how you arrived at your final design and implementation.  How did
the design evolve from the first version to the final version?} 

\plt{Don't just say what you changed, say why you changed it.  The needs of the
client should be part of the explanation.  For example, if you made changes in
response to usability testing, explain what the testing found and what changes
it led to.}

Over the course of the project, our design evolved significantly in response to feedback from various sources,
including peer reviews, TA and supervisor feedback, and usability testing.
At the beginning of the project, our planned flow was to use computer vision to read the game state off a
physical board, and then have a trained reinforcement learning model provide move recommendations
through an interface, acting as a sort of coach for the user playing the game.
Obviously, our design now has evolved significantly from this initial concept.

The first major change in our plans came quite early in the process.
We began by searching for existing Catan simulators, as designing a full sim for an RL model to
train on would have been a significant undertaking.
We came across Catanatron, an open-source Catan simulator written in Python, and decided to use it as
the core engine of our project.
Having made this decision, we began work on integrating with Catanatron, taking the time to understand
the codebase and design of the simulator, as well as its existing bots.
We then spent some time designing and implementing some simple bots to test our integration,
as well as incorporating Catanatron unit tests into our CI pipeline to ensure our integration didn't
break any existing functionality.
This decision to use Catanatron was a major turning point for our project, as it meant we could get
straight to work on our RL model.
We progressed with this model quite quickly early on, and by the time of our rev0 demonstration,
we had a working model that could play a game of Catan against the existing bots in Catanatron,
and beat a random bot consistently.
This was a significant milestone, and we decided as a team to defer the development of the
computer vision component indefinitely, in order to put more effort into the bot itself.
This meant we would no longer be reading off a physical board, and we moved away from the idea of a coach
to a more traditional bot that could be played against through Catanatron's CLI.

At this stage we put our full effort into improving our bot, implementing a variety of RL techniques
to improve training, such as reward shaping, curriculum learning, and self-play.
The bot was progressing steadily, and we were happy with the results we were seeing, but in a
meeting with our supervisor Prof. Istvan David, he expressed concern that while our work was impressive,
the only method we had of showing it was through Catanatron's CLI, and for a rev1 demo presented
to an audience less familar with RL, it would likely appear little progress had been made since rev0.
This was a very valid point, and we decided to pivot our design once again, this time towards the
development of a front-end interface that would allow users to play against our bot in a more
engaging way, as well as better demonstrate the progress we had made with the bot itself.
This was a significant change, as it meant we would be developing a full front-end interface,
which was a significant undertaking, but we felt it was necessary in order to properly showcase our work.
With only a couple weeks left until the rev1 demo, out first take on this UI was quite basic, but
it was a good starting point and went over quite well in the demo.

Now having a working UI, we realised this was the ideal situation for usability testing, so we
ran a round of interviews and surveys to gather feedback on the UI and overall user experience.
This exposed a few key areas for improvement in the design.
Similar to Prof. Smith in the demo, the users found the action log colours in the UI to be hard on the eyes,
especially the red and blue text on a black background.
They also disliked the look of the board, as some tile colours (such as lumber vs. wool) were hard to differentiate,
and the overall look was quite bland.
The action log colours were chosen to reflect the colour of each player, so we didn't remove them entirely,
but we played around with different hue and saturation levels to find a more visually appealing colour scheme.
We also redesigned the entire board, changing tile colours and textures.
Users also commented on other issues, such as the lack of a clear indication of which bot was which, as well
as what the current dice roll was, so we added a bot name label and a dice roll display to the UI.
Other changes were also made, you can see the details in the usability testing document in the extras folder.
Follow-up usability testing showed a noticeable increase in the UX metrics we were tracking, indicating our
changes were valuable.

While these changes were being made, we had continued meetings with Istvan.
In one of these meetings, he expressed that having the sole goal of creating the best bot was a bit narrow,
since likely someone will likely come along in the future with more compute resources and time and create
something better.
This was a good point, so we decided to take an idea given by Prof. Smith in the rev1 demo,
and additionally utilize our bot as a learning tool.
We accomplished this by adding an "explain move" feature to the UI,
where users can receive LLM-generated explanations of why a bot made the move it did, and allowing users
to walk back through a replay of the game.
This allowed users to develop strategic understanding from the bots that they could then apply in
their own games, supporting experiential learning.

\section{Design Decisions (LO12)}

\plt{Reflect and justify your design decisions.  How did limitations,
 assumptions, and constraints influence your decisions?  Discuss each of these
 separately.}



\section{Economic Considerations (LO23)}

\plt{Is there a market for your product? What would be involved in marketing your 
product? What is your estimate of the cost to produce a version that you could 
sell?  What would you charge for your product?  How many units would you have to 
sell to make money? If your product isn't something that would be sold, like an 
open source project, how would you go about attracting users?  How many potential 
users currently exist?}

Our repo is open source, and we have no plans to sell anything produced during this project.
Part of our motivation for this project was to push the field of Catan bots forward, and charging for
access to our bot/arena would go against this motivation.
As such, we'll focus instead in this section on how we will go about attracting users to our bot,
and how many potential users currently exist.

Multiple members of our team are active in the Catan community, and we've made connections with other
devs working on Catan bots throughout the project, thanks to are discussions in the Catanatron discord server.
This has given us a good network of potential dev users for our system, and we plan to invite them
to hook their own bots up to our bot arena, allowing them to benchmark their bots against ours.
We are already talking to a few devs with promising bots, and hope to have their bots integrated soon.
This will be a good way to attract users who are already interested in Catan bots, and we hope it will
give our platform a reputation as a hub for good Catan bots, attracting more devs.

Of course, our focus is not only on dev users, but with more discussion about our tool, we believe
we can attract a good number of regular Catan players as well, who are interested in playing against a strong bot
or using the bot as a learning tool to improve their own play.
We've already had some players from Colonist.io, a popular online Catan platform, express interest in
testing themselves against our bot, and we plan to reach out to more players in the community through
social media and Catan forums.
Should a bot hosted on our platform ever surpass the strength of the best human players, we believe this would
generate a lot of buzz and attract a large number of users to our platform, and this may not be all
that far-fetched of a possibility given the progress in the field of Catan bots in recent years.

Catan is the most popular modern board game in the world, and Colonist.io alone has 4.3 million registered users,
so while many may not be interested in testing themselves against or learning from our bot,
we still have a large potential user base.

\section{Reflection on Project Management (LO24)}

\plt{This question focuses on processes and tools used for project management.}

\subsection{How Does Your Project Management Compare to Your Development Plan}

\plt{Did you follow your Development plan, with respect to the team meeting plan, 
team communication plan, team member roles and workflow plan.  Did you use the 
technology you planned on using?}

\subsection{What Went Well?}

\plt{What went well for your project management in terms of processes and 
technology?}

\subsection{What Went Wrong?}

\plt{What went wrong in terms of processes and technology?}

\subsection{What Would you Do Differently Next Time?}

\plt{What will you do differently for your next project?}

\section{Ethics and Equity (LO18)}

\plt{Describe how your team applied the principle of equity and universal design
to ensure equitable treatment of all stakeholders.}

\section{Reflection on Capstone}

\plt{This question focuses on what you learned during the course of the capstone project.}

\subsection{Which Courses Were Relevant}

\plt{Which of the courses you have taken were relevant for the capstone project?}

\subsection{Knowledge/Skills Outside of Courses}

\plt{What skills/knowledge did you need to acquire for your capstone project
that was outside of the courses you took?}

\end{document}

